\section{Local limit theorem details}
\label{app:F}

Throughout, \(M_h(\theta)\) denotes the Fourier-tilted transfer matrix of the
synchronized product chain on \(E_h\times E_h\): for a product edge
\(E=\bigl((x,y)\to(x',y')\bigr)\) carrying the rank-two increment
\(g(E)=\eta(y\to y')-\eta(x'\to x)\) of Section~\ref{sec:fold}, we set
\(M_h(\theta)_{(x,y),(x',y')}=e^{i\langle\theta,g(E)\rangle}\) when \(E\) is an
edge and \(0\) otherwise. Thus \(M_h(0)\) is the adjacency matrix of the
product graph, which is primitive by Proposition~\ref{prop:graph}(5); write
\(\varrho=\rho(M_h(0))\).

\subsection{Wielandt equality and the lattice obstruction}

\begin{lemma}
\label{lem:F1}
For every \(\theta\in[-\pi,\pi)^2\) with \(\theta\neq0\) we have
\(\rho(M_h(\theta))<\varrho\).
\end{lemma}

\begin{proof}
Entrywise \(|M_h(\theta)|=M_h(0)\), so \(\rho(M_h(\theta))\le\varrho\).
Suppose equality holds. Since \(M_h(0)\) is primitive, Wielandt's theorem
yields \(\varphi\in\mathbb R\) and a diagonal unitary \(D\) with
\(M_h(\theta)=e^{i\varphi}DM_h(0)D^{-1}\); equivalently there is
\(\psi:E_h\times E_h\to\mathbb R/2\pi\mathbb Z\) with
\[
\langle\theta,g(E)\rangle\equiv\varphi+\psi(u')-\psi(u)\pmod{2\pi}
\]
for every product edge \(E=(u\to u')\). Summing along a closed product walk
\(W\) of length \(L\), the coboundary telescopes, leaving
\begin{equation}\label{eq:F1cocycle}
\Bigl\langle\theta,\ \textstyle\sum_{E\in W}g(E)\Bigr\rangle\equiv L\varphi
\pmod{2\pi}.
\end{equation}

Let \(c_1,c_2,c_3\) be the three equal-length closed walks of length \(L\) at
the common root \(z_0\) recorded for this \(h\) in Appendix~\ref{app:A}. For
each pair \((i,j)\) let \(W_{ij}\) be the closed product walk at
\((z_0,z_0)\) that traverses \(c_i\) backwards in the first coordinate and
\(c_j\) forwards in the second. Then \(W_{ij}\) has length \(L\) and, by the
definition of \(g\),
\[
\sum_{E\in W_{ij}}g(E)=\Psi(c_j)-\Psi(c_i),
\]
the projected Parikh difference. Taking \(i=j\) makes the left-hand side of
\eqref{eq:F1cocycle} zero, whence \(L\varphi\equiv0\pmod{2\pi}\); the constant
\(\varphi\) is thereby eliminated with no separate argument. Substituting this
back into \eqref{eq:F1cocycle} for general \(i,j\) gives
\[
\bigl\langle\theta,\ \Psi(c_j)-\Psi(c_i)\bigr\rangle\equiv0\pmod{2\pi}.
\]
In particular \(\langle\theta,d_1\rangle\equiv\langle\theta,d_2\rangle\equiv0\)
for \(d_1=\Psi(c_2)-\Psi(c_1)\) and \(d_2=\Psi(c_3)-\Psi(c_1)\).

This is the only point at which the covolume of the certificate is used.
Appendix~\ref{app:A} records \(\det[d_1\ d_2]=\pm1\), so \(\{d_1,d_2\}\) is a
\(\mathbb Z\)-basis of \(\mathbb Z^2\); had the determinant been some \(q\)
with \(|q|>1\), the conclusion below would only give
\(\theta\in(2\pi/q)\mathbb Z^2\) and nontrivial characters would survive.
Unimodularity gives \(\langle\theta,v\rangle\equiv0\pmod{2\pi}\) for every
\(v\in\mathbb Z^2\), hence for \(v=e_1,e_2\), so \(\theta\in2\pi\mathbb Z^2\),
i.e.\ \(\theta=0\) in \([-\pi,\pi)^2\) --- a contradiction.
\end{proof}

\subsection{Uniform spectral decay away from \texorpdfstring{\(\theta=0\)}{theta=0}}

\begin{lemma}
\label{lem:F2}
For every \(\delta>0\) there exist \(\lambda_\delta<\varrho\) and
\(Q<\infty\) with
\[
\|M_h(\theta)^N\|\le Q\,\lambda_\delta^{\,N}
\qquad\text{for all }N\ge0\text{ and all }\theta\in K_\delta
:=\{\delta\le\|\theta\|_\infty\le\pi\}.
\]
\end{lemma}

\begin{proof}
The entries of \(M_h(\theta)\) are trigonometric polynomials, so
\(\theta\mapsto M_h(\theta)\) is continuous; in fixed finite dimension the
roots of the characteristic polynomial depend continuously on the entries, so
\(\theta\mapsto\rho(M_h(\theta))\) is continuous. By Lemma~\ref{lem:F1} it is
strictly less than \(\varrho\) at every point of the compact set
\(K_\delta\), hence
\[
\lambda'_\delta:=\max_{\theta\in K_\delta}\rho(M_h(\theta))<\varrho .
\]
A pointwise Gelfand estimate would not suffice here, since its implied
constant depends on \(\theta\). Fix \(\lambda_\delta\) with
\(\lambda'_\delta<\lambda_\delta<\varrho\) and let
\(C_\delta=\{z:|z|=\lambda_\delta\}\). For
\((z,\theta)\in C_\delta\times K_\delta\) we have
\(z\notin\operatorname{spec}(M_h(\theta))\), so the resolvent
\(R(z,\theta)=(zI-M_h(\theta))^{-1}\) exists; matrix inversion is continuous
where defined, so \((z,\theta)\mapsto\|R(z,\theta)\|\) is continuous on the
compact set \(C_\delta\times K_\delta\) and
\[
Q_0:=\sup_{C_\delta\times K_\delta}\|R(z,\theta)\|<\infty .
\]
Since the whole spectrum of \(M_h(\theta)\) lies strictly inside
\(C_\delta\), the holomorphic functional calculus gives
\(M_h(\theta)^N=\frac{1}{2\pi i}\oint_{C_\delta}z^NR(z,\theta)\,dz\), whence
\[
\|M_h(\theta)^N\|\le\frac{1}{2\pi}\cdot2\pi\lambda_\delta\cdot
\lambda_\delta^{\,N}\,Q_0=Q\,\lambda_\delta^{\,N},
\qquad Q:=Q_0\lambda_\delta,
\]
uniformly on \(K_\delta\).
\end{proof}

\subsection{Uniformity in endpoints, bounded displacement and lengths \texorpdfstring{\(N+r\)}{N+r}}

\begin{proposition}
\label{prop:F3}
Fix \(r_0\in\mathbb N\) and a finite \(V\subset\mathbb Z^2\). Let
\(\mathcal N_N(\alpha,\omega;v)\) be the number of length-\(N\) product paths
from \(\alpha\) to \(\omega\) with \(\sum_E g(E)=v\), and let
\(\mathcal W_N(\alpha,\omega)=(M_h(0)^N)_{\alpha,\omega}\). Then
\[
\mathcal N_{N+r}(\alpha,\omega;v)
=\mathcal W_{N+r}(\alpha,\omega)\,
\frac{1}{2\pi N\sqrt{\det\Sigma_{D,h}}}\,\bigl(1+o(1)\bigr),
\]
where the \(o(1)\) is uniform simultaneously over all state pairs
\((\alpha,\omega)\), all \(v\in V\), and all integers \(|r|\le r_0\).
\end{proposition}

\begin{proof}
Fourier inversion on \(\mathbb Z^2\) is exact:
\[
\mathcal N_{N+r}(\alpha,\omega;v)=\frac1{(2\pi)^2}
\int_{[-\pi,\pi]^2}e^{-i\langle\theta,v\rangle}
\bigl(M_h(\theta)^{N+r}\bigr)_{\alpha,\omega}\,d\theta .
\]
Split at \(\|\theta\|_\infty=\delta\). On \(K_\delta\), Lemma~\ref{lem:F2}
bounds the integrand by \(Q\lambda_\delta^{N+r}\) with
\(\lambda_\delta<\varrho\); this bound involves neither \(\alpha,\omega\) nor
\(v\), so that contribution is \(O(\lambda_\delta^{N})=o(\varrho^{N}/N)\)
uniformly in all three parameters.

Near the origin, \(M_h(0)\) is primitive, so \(\varrho\) is a simple isolated
eigenvalue. By analytic perturbation theory there are \(\delta_0>0\), an
analytic simple eigenvalue \(\varrho(\theta)\) and an analytic rank-one
spectral projection \(\Pi(\theta)\) on \(\|\theta\|_\infty<\delta_0\), with the
remaining spectrum inside \(\{|z|\le\varrho''\}\) for some
\(\varrho''<\varrho\); hence
\(M_h(\theta)^{N+r}=\varrho(\theta)^{N+r}\Pi(\theta)+O((\varrho'')^{N+r})\)
uniformly. Write \(\varrho(\theta)=\varrho\,e^{p(\theta)}\) with \(p\)
analytic and \(p(0)=0\). By Lemma~\ref{lem:cov2} and the \(S_3\) letter
frequencies \((1/3,1/3,1/3)\) the drift vanishes, so \(\nabla p(0)=0\) and
\(p(\theta)=-\tfrac12\theta^{\!\top}\Sigma_{D,h}\theta+O(\|\theta\|^3)\).

Substituting \(\theta=\phi/\sqrt N\) and using \(\nabla p(0)=0\),
\[
\varrho(\theta)^{N}=\varrho^{N}
\exp\bigl(-\tfrac12\phi^{\!\top}\Sigma_{D,h}\phi\bigr)(1+o(1)),
\]
and \(\Pi(\theta)\to\Pi(0)\). Gaussian integration produces the factor
\((2\pi N)^{-1}(\det\Sigma_{D,h})^{-1/2}\).

The three uniformities are now read off. \emph{Endpoints:} there are finitely
many pairs \((\alpha,\omega)\) and the dependence enters only through the entry
\(\Pi(\theta)_{\alpha,\omega}\), so the maximum of finitely many \(o(1)\)
terms is \(o(1)\). \emph{Displacement:} \(v\) enters only through
\(e^{-i\langle\theta,v\rangle}=e^{-i\langle\phi,v\rangle/\sqrt N}\), and
\(\bigl|e^{-i\langle\phi,v\rangle/\sqrt N}-1\bigr|\le\|\phi\|\,\|v\|/\sqrt N\),
which is uniform over the fixed finite set \(V\). \emph{Length:}
\(\varrho(\theta)^{N+r}=\varrho(\theta)^{N}\varrho(\theta)^{r}\) and
\(\varrho(\theta)^{r}\to\varrho^{r}\) uniformly for \(|r|\le r_0\), while
\((2\pi(N+r))^{-1}=(2\pi N)^{-1}(1+O(1/N))\).
\end{proof}

\begin{remark}
This is exactly the form consumed by Lemma~\ref{lem:boundary}. After peeling,
at most \(D_h\) edges are removed at each end, so the surviving core portion
has length \(N-r\) with \(0\le r\le 2D_h\); its endpoints are the finitely many
core states at which the path enters and leaves \(E_h\); and by
Lemma~\ref{lem:fold} the required displacement is \(\beta\) adjusted by the
Parikh contribution of the peeled edges, so it lies in the fixed finite set
\(\{v:\|v\|_\infty\le m+2D_h\}\). The resulting relative error is therefore a
single \(o(1)\) independent of the term, and since every term of the identity
of Lemma~\ref{lem:pathid} is non-negative, summation preserves it.
\end{remark}
